\chapter{Đi và Về}
\markboth{Đi và Về}{Going and Returning}

\section{Đi xa để thấy quê nhà không hề nhỏ.}
\begin{truyen}{Đi xa để thấy quê nhà không hề nhỏ.}{Distance can deepen belonging.}
\chuhoa{R}{a thành phố học, em từng thấy quê mình chẳng có gì ghê gớm. Ở lâu rồi mới thấm: bữa cơm nhà, con đường cũ, tiếng người quen gọi nhau mới là thứ đỡ mình sau những ngày bơ phờ. Đi xa không phải để chê chỗ cũ; nhiều khi là để biết mình đã được nuôi bằng thứ gì.}
\textit{(Leaving for the city, the student once thought home was too ordinary. Yet after months of exhaustion in unfamiliar rhythms, old roads, family meals, familiar voices, and overlooked peace all returned with new meaning. Some journeys do not make us despise where we came from; they teach us to honor what raised us.)}
\end{truyen}
\begin{baihoc}
Đi xa một đoạn mới biết cái mình từng xem thường lại là chỗ đỡ lưng cho mình.
\textit{(Going far is not always leaving behind; often it is the way to understand more deeply what we belong to.)}
\end{baihoc}
\begin{ghinhoanh}
\textbf{Short English to remember:}

Going far can deepen home.
\end{ghinhoanh}
\begin{tuhocnhanh}
1. Điều gì ở nơi mình thuộc về mà em chỉ thấy quý khi đi xa? \textit{(What about your place of belonging do you only value when you are far away?)}

2. Em có đang xem nhẹ điều gì rất nền tảng trong đời mình không? \textit{(Are you underestimating something foundational in your life?)}
\end{tuhocnhanh}
\ngancach

\section{Đi nhiều nhưng không biết mình đang về đâu.}
\begin{truyen}{Đi nhiều nhưng không biết mình đang về đâu.}{Movement is not the same as direction.}
\chuhoa{A}{nh lúc nào cũng bận, lịch kín, điện thoại reo liên tục, nhìn ngoài vào rất có vẻ đang tiến. Nhưng nhiều đêm nằm xuống, chính anh cũng không biết rốt cuộc mình đang chạy vì cái gì. Đi như bị đuổi mà không biết đích ở đâu thì cũng chỉ là mệt thêm.}
\textit{(He was constantly busy, constantly occupied, constantly moving. To outsiders he looked energetic and fast-moving. Only he knew that some nights, lying down, he could not tell what destination his life was actually heading toward.)}
\end{truyen}
\begin{baihoc}
Chạy hùng hục mà sai hướng thì cũng chỉ là tự bào mình nhanh hơn.
\textit{(Moving a lot does not replace the question of life direction.)}
\end{baihoc}
\begin{ghinhoanh}
\textbf{Short English to remember:}

Movement is not direction.
\end{ghinhoanh}
\begin{tuhocnhanh}
1. Em đang rất bận hay đang thật sự có hướng? \textit{(Are you merely busy, or do you truly have direction?)}

2. Nếu dừng lại một lúc, em có trả lời được mình đang đi vì điều gì không? \textit{(If you paused, could you answer why you are going the way you are?)}
\end{tuhocnhanh}
\ngancach

\section{Đi để học, về để trao lại.}
\begin{truyen}{Đi để học, về để trao lại.}{A good journey returns with service.}
\chuhoa{M}{ột người đi xa học được nhiều cái hay nhưng không mang về để vênh mặt với quê mình. Anh chọn lọc cái dùng được, sửa cho hợp, rồi làm từng chút để nơi cũ đỡ thiếu hơn. Đi cho đã rồi quay về chê bai thì dễ; mang được chút ích lợi về mới khó.}
\textit{(A man learned many things elsewhere but did not use them to despise his hometown. When he returned, he shared what was useful, adapted it wisely, and helped improve the old place. A beautiful journey not only enlarges the traveler, but also warms the place of return.)}
\end{truyen}
\begin{baihoc}
Đi xa mà còn nhớ mang cái tốt về trả lại mới là đi không phí.
\textit{(To go without forgetting the responsibility to return is a worthy kind of maturity.)}
\end{baihoc}
\begin{ghinhoanh}
\textbf{Short English to remember:}

Learn far, give back home.
\end{ghinhoanh}
\begin{tuhocnhanh}
1. Điều gì em muốn mang về cho nơi mình xuất phát? \textit{(What do you want to bring back to the place you came from?)}

2. Em học để hơn người hay để có ích hơn? \textit{(Are you learning to be above others or to be more useful?)}
\end{tuhocnhanh}
\ngancach

\section{Biết quay về với gia đình sau những tham vọng dài.}
\begin{truyen}{Biết quay về với gia đình sau những tham vọng dài.}{Returning is sometimes the bravest correction.}
\chuhoa{S}{au nhiều năm chỉ biết cắm đầu vào việc, anh mới giật mình vì mình bỏ lỡ quá nhiều bữa cơm, quá nhiều lần con đợi, quá nhiều mùa cha mẹ già thêm. Về sớm một bữa, ngồi ăn cho ra bữa, hỏi nhau mấy câu tử tế, nghe thì nhỏ nhưng với anh lúc đó lại là chuyện lớn. Có khi quay về nhà mới là cú rẽ khó nhất.}
\textit{(After many years of chasing work, he realized how many family dinners, waiting children, and aging parents he had missed. Returning earlier, sitting for a proper meal, asking more carefully about others' lives, no longer felt small. Some returns are major decisions that rescue the parts of life wearing thin.)}
\end{truyen}
\begin{baihoc}
Còn biết quay về nghĩa là vẫn chưa lạc tới mức hết đường cứu.
\textit{(Knowing how to return to what is essential is a sign that a person has not fully lost themselves.)}
\end{baihoc}
\begin{ghinhoanh}
\textbf{Short English to remember:}

Returning can rescue what matters.
\end{ghinhoanh}
\begin{tuhocnhanh}
1. Điều cốt lõi nào trong đời em đang cần được quay về? \textit{(What essential part of your life needs returning to?)}

2. Em đã xa điều quan trọng ấy từ lúc nào? \textit{(Since when have you been drifting away from it?)}
\end{tuhocnhanh}
\ngancach

\section{Đi tiếp nhưng không quên trở về với chính mình.}
\begin{truyen}{Đi tiếp nhưng không quên trở về với chính mình.}{The longest journey still needs inner return.}
\chuhoa{T}{rưởng thành đâu có nghĩa là đứng im. Vẫn phải đi, vẫn va chạm, vẫn đổi nghề, đổi chỗ, đổi suy nghĩ. Nhưng đi kiểu gì thì đi, lâu lâu vẫn phải quay vào trong mà hỏi: mình đang sống thật không, hay chỉ đang diễn; mình đang tử tế không, hay đang khô dần. Còn hỏi được mấy câu đó thì chưa lạc hẳn.}
\textit{(Maturity does not mean standing still. We still go, learn, collide, and change. But if through all those journeys we still know how to return inward and ask whether we are living honestly, living decently, and being led by the right things, then we are not entirely lost.)}
\end{truyen}
\begin{baihoc}
Đi và về không chỉ là chuyện đường xa. Nó là chuyện mình có còn giữ được gốc khi đời kéo mình đi đủ hướng hay không.
\textit{(Going and returning is not only about distance; it is the rhythm of a person who opens outward while preserving the root.)}
\end{baihoc}
\begin{ghinhoanh}
\textbf{Short English to remember:}

Go far, return inward.
\end{ghinhoanh}
\begin{tuhocnhanh}
1. Em đang đi mạnh ở bên ngoài, nhưng bên trong có còn chỗ để trở về không? \textit{(You may be moving strongly outside, but do you still have an inner place to return to?)}

2. Câu hỏi nào nên trở thành điểm quay về thường xuyên của em? \textit{(Which question should become your regular point of return?)}
\end{tuhocnhanh}
\ngancach
